\documentclass[11pt,a4paper]{article}

\usepackage[margin=2.4cm]{geometry}
\usepackage{amsmath,amssymb}
\usepackage{booktabs}
\usepackage{longtable}
\usepackage{array}
\usepackage{xcolor}
\usepackage[colorlinks=true,linkcolor=blue!50!black,urlcolor=blue!55!black,citecolor=blue!50!black]{hyperref}
\usepackage{fancyvrb}
\usepackage{tcolorbox}
\usepackage{enumitem}
\usepackage{titlesec}

\titleformat{\section}{\normalfont\Large\bfseries\color{blue!45!black}}{\thesection}{1em}{}
\titleformat{\subsection}{\normalfont\large\bfseries\color{blue!30!black}}{\thesubsection}{1em}{}

\newcommand{\code}[1]{\texttt{\small #1}}
\newcommand{\Rey}{\mathrm{Re}}

\title{\bfseries State of the Art:\\ Newton Globalization and Solver Performance\\ for Implicit CFD}
\author{Survey for the SFEM CVFEM semi-structured multigrid solver}
\date{\today}

\begin{document}
\maketitle

\begin{tcolorbox}[colback=red!4,colframe=red!55!black,title=\bfseries Headline finding]
Our solver needs $\sim$\textbf{1000} Krylov iterations per Newton step and fails above
$\Rey \approx 200$. Published incompressible Navier--Stokes solvers using an
augmented-Lagrangian preconditioner need \textbf{4--8} iterations per Newton step and reach
$\Rey = 10\,000$. Plain Newton with \emph{no} globalization at all is documented to reach
$\Rey \approx 2500$.

Meanwhile our \emph{cost per cycle} is normal: our V-cycle runs at 6.06\,M nodes/s, and FUN3D
on the same Grace silicon completes a full nonlinear iteration at $\approx 6.0$\,M points/s.

\medskip
\textbf{The conclusion is therefore that the linear solve, and specifically the smoother, is
the binding constraint --- not the continuation strategy.}
\end{tcolorbox}

\section{Scope and baseline}

This surveys how established implicit CFD codes get Newton to converge at high Reynolds
number, and what throughput they achieve, in order to judge our own numbers. Our reference
configuration throughout:

\begin{center}
\begin{tabular}{@{}ll@{}}
\toprule
Discretisation & colocated FV/CVFEM, hexahedra, Rhie--Chow stabilisation \\
Solver & Newton + geometric multigrid V-cycle, block-Jacobi smoother, block size 4 \\
Hardware & one NVIDIA Grace socket, 72 Neoverse V2 cores \\
Measured & 242\,500 dof, one V-cycle $\approx$ 10\,ms; 646\,MDOF/s operator apply \\
Converged run & $\Rey=100$, 33\,124 dof: 32 Newton steps, $\approx$30\,000 linear iterations \\
Failure & diverges above $\Rey \approx 200$--$400$ \\
\bottomrule
\end{tabular}
\end{center}

\section{Pseudo-transient continuation: the formulas}

The underlying iteration adds a pseudo-time term to the Jacobian, recovering Newton as
$\Delta t \to \infty$:
\begin{equation}
\left[ \frac{M}{\Delta t} + \frac{\partial R}{\partial U} \right] \Delta U^k = -R(U^k),
\qquad U^{k+1} = U^k + \omega^k \Delta U^k,
\qquad \mathrm{CFL} = \frac{\lambda_{\max}\,\Delta t_\kappa}{L_\kappa}.
\end{equation}
Convergence theory: Kelley \& Keyes, \emph{SINUM} \textbf{35}(2), 1998, 508--523,
\href{https://doi.org/10.1137/S0036142996304796}{doi:10.1137/S0036142996304796}.

\subsection{SER --- switched evolution relaxation}
Mulder \& van Leer, \emph{JCP} \textbf{59} (1985) 232--246
(\href{https://homepage.tudelft.nl/a1k53/Papers/Journal/JCP59_Mulder.pdf}{PDF}). The CFL is
scaled by the inverse residual ratio:
\begin{equation}
\mathrm{CFL}^k = \min\!\left( \mathrm{CFL}^{k-1}\,
\frac{\lVert R^{k-1}\rVert_2}{\lVert R^{k}\rVert_2},\; \mathrm{CFL}_{\max} \right)
\end{equation}
A variant references the \emph{initial} residual instead of the previous one:
\begin{equation}
\mathrm{CFL}^n = \frac{\mathrm{CFL}_0}{r_n}, \qquad
r_n = \max\!\left( \frac{\lVert R(q^n)\rVert_2}{\lVert R(q^0)\rVert_2},\;
\frac{\lVert R(q^n)\rVert_\infty}{\lVert R(q^0)\rVert_\infty} \right)
\end{equation}
PETSc's \href{https://petsc.org/main/manualpages/TS/TSPSEUDO/}{\code{TSPSEUDO}} implements both
forms, with growth factor \code{-ts\_pseudo\_increment} and cap \code{-ts\_pseudo\_max\_dt}.

\subsection{EXPur --- driven by the under-relaxation factor}
Ceze \& Fidkowski, AIAA 2013-2686
(\href{https://public.websites.umich.edu/~kfid/MYPUBS/Ceze_Fidkowski_2013_PTC.pdf}{PDF}), eqn.~(13):
\begin{equation}
\mathrm{CFL}^{k+1} =
\begin{cases}
\beta\,\mathrm{CFL}^k & \omega^k = 1 \quad (\beta>1)\\
\mathrm{CFL}^k & \omega_{\min} < \omega^k < 1\\
\kappa\,\mathrm{CFL}^k & \omega^k < \omega_{\min} \quad (\kappa<1)
\end{cases}
\end{equation}
Parameters used: $\omega_{\min}=0.01$, $\beta = 1.05$--$2.0$, $\kappa = 0.1$.

\subsection{RDM --- residual difference method (bounded SER)}
Same reference, eqn.~(14). Growth is capped at $\beta$, with equality only when the residual
vanishes:
\begin{equation}
\mathrm{CFL}^k = \min\!\left( \mathrm{CFL}^{k-1}
\beta^{\left[(\lVert R^{k-1}\rVert_2 - \lVert R^{k}\rVert_2)/\lVert R^{k-1}\rVert_2\right]},\;
\mathrm{CFL}_{\max}\right)
\end{equation}
The \emph{monotonic} variant mRDM forces the exponent to zero when the residual increases.

\section{What the production codes actually do}

\subsection{FUN3D: HANIM}
Wang, Diskin, Nielsen \& Liu, AIAA 2021-0857
(\href{https://ntrs.nasa.gov/api/citations/20205010283/downloads/fun3d_hanim.pdf}{PDF}). The
legacy scheme was a \emph{linear} CFL ramp specified as two namelist arrays (CFL $1 \to 500$
over 200 iterations) with 15 multicolour Gauss--Seidel sweeps and no Krylov method.

HANIM replaced it with a module hierarchy --- preconditioner, GCR, realizability check,
nonlinear control, CFL adaption --- in which \textbf{any module reporting failure discards the
update, cuts the CFL, and repeats from the same state}:
\begin{center}
\begin{tabular}{@{}ll@{}}
\toprule
success & $\mathrm{CFL} \leftarrow \mathrm{CFL} \times 2$ \\
failure & $\mathrm{CFL} \leftarrow \mathrm{CFL} \times 0.1$ \\
bounds  & $\mathrm{CFL}_{\min} = 10^{-12}$, $\mathrm{CFL}_{\max} = 10^{10}$ \quad (22 orders) \\
\bottomrule
\end{tabular}
\end{center}
On failure only the \emph{Jacobian diagonal} is updated for the new CFL; fluxes are not
re-linearized. The Krylov success criterion is deliberately loose,
$\lVert r_k\rVert/\lVert r_0\rVert \le \mu_{\mathrm{GCR}} = 0.92$ --- an 8\,\% reduction per
outer step. The line search is a \emph{quadratic search with a closed-form minimiser} that
reuses the linear residual GCR has already computed, so it costs one extra residual
evaluation rather than a sequence:
\begin{equation}
\omega = -\frac{g(1)-f(0)}{2\,(f(1)-g(1))},\qquad
g(1) = \left\lVert \tfrac{V}{\Delta\tau}\Delta Q + \tfrac{\partial R}{\partial Q}\Delta Q + R(Q^n)\right\rVert
\end{equation}

\begin{center}
\begin{tabular}{@{}lrrrr@{}}
\toprule
Case & Grid & Cores & Iterations (legacy $\to$ HANIM) & Speedup \\
\midrule
Hemisphere-cylinder & 9.0\,M & 160 & $8976 \to 675$ & 2.6$\times$ \\
Hemisphere-cylinder & 71.4\,M & 1400 & $10\,907 \to 1874$ & 2.1$\times$ \\
Supersonic duct & 3.17\,M & 240 & $9721 \to 1283$ & 2.2$\times$ \\
Juncture flow & 13.0\,M & 320 & no convergence in 72\,h $\to$ converged & 6.8$\times$ \\
Juncture flow & 37.1\,M & 756 & --- & 7.8$\times$ \\
\bottomrule
\end{tabular}
\end{center}

\subsection{NASA LAURA / HyperSolve: the most directly implementable controller}
Thompson \& O'Connell
(\href{https://ntrs.nasa.gov/api/citations/20200002615/downloads/20200002615.pdf}{NASA PDF}),
Algorithm 1, verbatim:
\begin{Verbatim}[fontsize=\small,frame=single,framesep=3mm]
CFL <- 1.0 ;  Q_safe <- Q
while |R(Q)|_2 > tolerance:
    Dt_i  <- CFL * V_i / sum_f (|u_f| + c) A_f
    dQ    <- (M/Dt + dR/dQ)^-1 R_t
    eta   <- line search s.t. |R_t(Q)|_2 > |R_t(Q + eta dQ)|_2
    if eta > 0.1:  CFL <- 1.5*CFL ;  Q <- Q + eta dQ ;  Q_safe <- Q
    else:          CFL <- 0.1*CFL ;  Q <- Q_safe
\end{Verbatim}
The critical detail is that the time term appears on \emph{both} sides, so the merit function
and the Jacobian stay consistent:
$R_t = R + (M/\Delta t)\Delta Q$ and
$[M/\Delta t + \partial R/\partial Q]\Delta Q = -(M/\Delta t)\Delta Q - R(Q^n)$.

\subsection{SU2: the shipping controller}
From \code{config\_template.cfg}: \code{CFL\_ADAPT\_PARAM = (0.1, 2.0, 10.0, 1e10, 0.001, 0)},
i.e.\ factor-down, factor-up, min, max, \emph{acceptable linear-solver convergence}, start
iteration. The fifth parameter matters for us: \textbf{SU2 cuts the CFL when the linear solver
fails to converge sufficiently}. A naive SER controller ramps on residual history alone and
will happily accelerate while the linear solve silently degrades. Guidance: Euler and laminar
tolerate doubling to a very high cap; RANS wants factor-up $\approx 1.2$ and
$\mathrm{CFL}_{\max}\approx 10^3$.

\section{Reynolds continuation versus pseudo-transient}

The split is by regime, not preference. Compressible aerodynamic codes (FUN3D, SU2, CFL3D,
OVERFLOW, LAURA) use PTC exclusively: there is no natural continuation parameter, and
pseudo-time is the physically meaningful homotopy. The \emph{finite-element incompressible}
community uses Reynolds continuation, because $\nu$ is an explicit coefficient and the Stokes
problem at $\Rey \to 0$ is trivially solvable and supplies a free starting point.

\begin{tcolorbox}[colback=blue!4,colframe=blue!50!black,title=\bfseries The reference for our regime]
Farrell, Mitchell \& Wechsung, \emph{SISC} \textbf{41}(5), A3073--A3096
(\href{https://arxiv.org/pdf/1810.03315}{arXiv:1810.03315}). Augmented-Lagrangian
preconditioner for 3D stationary incompressible Navier--Stokes at high Reynolds number.

\medskip
Their schedule: solve \textbf{Stokes} first, then $\Rey = 10 \to 100$, then in steps of 100 up
to $\Rey = 10\,000$, each solution seeding the next. Newton globalized with a PETSc $L^2$
line search.
\end{tcolorbox}

\subsection{Their iteration counts --- the benchmark we should measure against}
Krylov iterations per Newton step, 3D lid-driven cavity:
\begin{center}
\begin{tabular}{@{}lrrrrrr@{}}
\toprule
Refinements & dof & $\Rey{=}10$ & 100 & 1000 & 5000 & 10\,000 \\
\midrule
1 & $1.0\times10^4$ & 2.50 & 4.33 & 6.00 & 8.00 & 14.00 \\
4 & $6.6\times10^5$ & 2.50 & 2.67 & 5.00 & 8.00 & 8.50 \\
\bottomrule
\end{tabular}
\end{center}
For contrast, PCD/LSC preconditioning (IFISS) needs 22--26 at $\Rey{=}10$ and 100--210 at
$\Rey{=}1000$. Against SIMPLE on the same problem, serial:
\begin{center}
\begin{tabular}{@{}lrrrr@{}}
\toprule
$\Rey$ & AL iters & AL time & SIMPLE iters & SIMPLE time \\
\midrule
10 & 4 & 0.21\,min & 515 & 1.06\,min \\
100 & 8 & 0.38 & 979 & 2.01 \\
200 & 10 & 0.46 & 1185 & 2.48 \\
\bottomrule
\end{tabular}
\end{center}
Each SIMPLE iteration is 22--26$\times$ cheaper and it still loses.
\textbf{Our $\sim$1000 cycles per Newton step sits in the SIMPLE column.}

\subsection{The alternative that needs no continuation at all}
Rebholz et al.\ (\href{https://arxiv.org/pdf/2402.12304}{arXiv:2402.12304}) compose one Picard
step with one Newton step. The Picard half supplies the stability Newton lacks --- global
stability for plain Newton does \emph{not} hold --- while retaining local quadratic
convergence.
\begin{center}
\begin{tabular}{@{}lr@{}}
\toprule
Method (same mesh, zero initial guess, no continuation) & Converges to \\
\midrule
Newton & $\Rey \lesssim 2500$ \\
Newton + line search & $\Rey \lesssim 3000$ \\
Picard & $\Rey \lesssim 4000$ \\
\textbf{Picard--Newton} & $\mathbf{\Rey \approx 10\,000\text{--}12\,000}$ \\
\bottomrule
\end{tabular}
\end{center}
\textbf{Plain Newton's own documented limit is $\Rey\approx2500$, an order of magnitude above
our $\Rey\approx200$--400.} This is the strongest single piece of evidence that our problem is
the linear solve rather than the globalization.

\section{Inexact Newton: Eisenstat--Walker forcing terms}
Eisenstat \& Walker, \emph{SISC} \textbf{17} (1996) 16--32
(\href{https://users.wpi.edu/~walker/Papers/forcing_terms,SISC_17,1996,16-32.pdf}{PDF}). This
directly attacks over-solving the linear system early in the Newton sequence. PETSc's default
(version 2):
\begin{equation}
\eta_k = \gamma \left( \frac{\lVert F_k\rVert}{\lVert F_{k-1}\rVert} \right)^{\alpha},
\qquad \alpha = \tfrac{1+\sqrt 5}{2} \approx 1.618,\;\; \gamma = 1,
\end{equation}
with $\eta_0 = 0.3$, $\eta_{\max} = 0.9$, and a safeguard
$\eta_k \leftarrow \max(\eta_k, \gamma\,\eta_{k-1}^{\alpha})$ when that exceeds $0.1$.

\section{Smoother strength: local Fourier analysis for Stokes}
Claus \& Bolten (\href{https://arxiv.org/pdf/2008.08719}{arXiv:2008.08719}), asymptotic
convergence factors:
\begin{center}
\begin{tabular}{@{}lcc@{}}
\toprule
Smoother & $\rho$ (periodic) & $\rho$ (Dirichlet) \\
\midrule
Vanka & \textbf{0.08} & 0.10 \\
Triad Gauss--Seidel & 0.26 & 0.36 \\
Triad Jacobi & 0.49 & \emph{divergent} \\
\bottomrule
\end{tabular}
\end{center}
Our point-block Jacobi is weaker than triad-Jacobi, which that study reports degrades to a
\emph{divergent} method under Dirichlet boundary conditions --- the case we run.

\section{Throughput: are our per-cycle costs reasonable?}

\subsection{FUN3D}
\begin{center}
\begin{tabular}{@{}llll@{}}
\toprule
Metric & Value & Hardware & Problem \\
\midrule
Full nonlinear iteration & $\approx$6.0\,M points/s & \textbf{GH200 Grace, 72 cores} & 5\,M points \\
3000 iterations & 2818\,s / 1639\,s / 430\,s & SKL node / EPYC node / A100 & 3.68\,M vertices \\
Derived & 3.91 / 6.73 / 25.6\,M points/s & SKL / EPYC / A100 & same \\
15 multicolour GS sweeps & 166\,ms $\to$ 30.9\,ms & SKL $\to$ A100 (81\,\% of peak BW) & 1.12\,M vertices \\
\bottomrule
\end{tabular}
\end{center}
\textbf{The comparison.} Our V-cycle is $242\,500\,\mathrm{dof} / 10\,\mathrm{ms} =
6.06$\,M~nodes/s. FUN3D on the same Grace silicon completes an entire nonlinear iteration
(residual, Jacobian, 15 Gauss--Seidel sweeps) at $\approx 6.0$\,M~points/s. Our V-cycle costs
about what a full production nonlinear iteration costs. NASA's own framing: \emph{``FUN3D is
memory bound: in general, performance scales with memory bandwidth.''}

\subsection{OpenFOAM: the throughput-versus-loading curve}
Galeazzo et al.\ 2024, one HLRS Hawk node ($2\times$EPYC 7742, 128 cores)
(\href{https://iopscience.iop.org/article/10.1088/1757-899X/1312/1/012009}{IOP MSE 1312 012009}):
\begin{center}
\begin{tabular}{@{}rrrrr@{}}
\toprule
Cells/rank & kcells/s/core & L3 miss:hit & MPI ratio & Parallel eff. \\
\midrule
210\,938 & 24.4 & 0.823 & 0.056 & 0.96 \\
62\,500 & 39.1 & 0.520 & 0.089 & 0.92 \\
26\,367 & 100.8 & 0.083 & 0.279 & 0.77 \\
\textbf{15\,259} & \textbf{106.0 (peak)} & 0.054 & 0.353 & 0.71 \\
7\,813 & 104.7 & 0.050 & 0.469 & 0.65 \\
3\,296 & 82.0 & 0.095 & 0.625 & 0.54 \\
122 & 7.0 & 0.342 & 0.875 & 0.19 \\
\bottomrule
\end{tabular}
\end{center}
Two consequences for us. First, our 842 nodes/core sits on the \emph{steep left flank},
between the 3296 and 122 rows, where MPI and parallel overhead dominate --- so our measured
throughput is more likely \textbf{understated} than inflated, and re-measuring at
$\approx$15\,000 nodes/core ($\approx$4.3\,M dof on 72 cores) should \emph{raise} it. Second,
peak throughput and peak parallel efficiency move in opposite directions: the peak sits at
71\,\% efficiency, not 95\,\%. Both numbers should be reported.

Our V-cycle costs 84\,knodes/s/core, inside OpenFOAM's 24--106\,kcells/s/core band for a
\emph{complete segregated timestep}. Two independent codes therefore agree that our per-cycle
cost is production-normal.

\section{Ranked recommendations}
\begin{enumerate}[leftmargin=*]
\item \textbf{Picard$\to$Newton composition.} We already have both linearizations: Newton's
Jacobian minus the $b^*(u^{k+1},\hat u,v)$ term \emph{is} Picard. Largest robustness gain per
line changed; reaches $\Rey\approx10\,000$ with no continuation at all.
\item \textbf{Eisenstat--Walker forcing terms.} A few lines; directly attacks over-solving.
\item \textbf{Stokes start plus $\Rey$ continuation} on Farrell's schedule, replacing our
current weak two-stage version.
\item \textbf{Full multigrid / nested iteration.} We already own the hierarchy; near-free.
\item \textbf{PTC with the LAURA controller} ($\approx$30 lines, fully specified above,
including the \code{Q\_safe} rollback that makes it safe).
\item \textbf{Per-step physicality limiting} ($\eta_{\max} = 10\,\%$ fractional change).
\item \textbf{HANIM in full}, if 1--6 prove insufficient.
\end{enumerate}

\begin{tcolorbox}[colback=orange!6,colframe=orange!60!black,title=\bfseries Caveat on the ordering]
None of items 1--6 will rescue a solver that needs 1000 V-cycles per Newton step. The
smoother must be measured alone --- as a stationary iteration run to its asymptote, with no
coarse space --- before any of this is worth implementing.
\end{tcolorbox}

\section{Where comparison is unsound}
\begin{itemize}[leftmargin=*]
\item FUN3D's published numbers are \textbf{compressible}, block size 5. Its incompressible
path (block size 4) has no published performance figures.
\item Farrell et al.\ use a \textbf{finite-element} discretisation with an augmented-Lagrangian
preconditioner; we use colocated FV with Rhie--Chow. The iteration counts are comparable as a
target, the preconditioner is not directly transplantable.
\item OpenFOAM figures are per \emph{segregated timestep}, not per Newton step.
\item The OHC-1 benchmark paper reports mutually inconsistent absolute MFVOPS figures
(\S4.2.2 versus Table 3); only its \emph{ratios} should be cited.
\item The SU2 smoother details (BCSR $5\times5$; ILU(0) tolerating roughly half the CFL of a
GMRES smoother) were \textbf{not verified against the primary text} and should be confirmed
before being designed around.
\end{itemize}

\end{document}
